\title{Direct Preference Optimization: Your Language Model is Secretly a Reward Model}

\begin{abstract}
Although large-scale unsupervised language models (LMs) acquire extensive world knowledge and some reasoning capabilities, precisely controlling their behavior remains challenging because their training is entirely unsupervised. Current approaches to achieve such control gather human annotations on the relative quality of model outputs and subsequently fine-tune the unsupervised LM to align with these preferences, often using reinforcement learning from human feedback (RLHF). Nevertheless, RLHF is a complicated and frequently unstable process that first involves fitting a reward model to capture human preferences, followed by fine-tuning the large unsupervised LM via reinforcement learning to maximize this learned reward while preventing excessive divergence from the original model. In this work, we propose a novel parameterization of the reward model used in RLHF that allows for direct extraction of the optimal policy in closed form, enabling us to address the standard RLHF problem using only a straightforward classification loss. This method, which we term \textit{Direct Preference Optimization} (DPO), is stable, effective, and computationally efficient, removing the necessity for sampling from the LM during fine-tuning or extensive hyperparameter optimization. Our experiments demonstrate that DPO can fine-tune LMs to align with human preferences at least as well as existing techniques. Importantly, fine-tuning with DPO surpasses PPO-based RLHF in controlling the sentiment of generated outputs and matches or enhances response quality in summarization and single-turn dialogue tasks, while being considerably simpler to implement and train.
\end{abstract}

\section{Introduction}
Large-scale unsupervised language models (LMs) trained on extensive datasets develop remarkable capabilities~\citep{chowdhery2022palm, brown2020language, touvron2023llama,bubeck2023sparks}. Nevertheless, these models learn from data created by humans who have diverse objectives, priorities, and expertise. Some of these objectives and skills may not be ideal to replicate; for instance, while it is beneficial for an AI coding assistant to \textit{recognize} common programming errors to correct them, when generating code, we want to steer the model towards the (often rare) high-quality coding skills embedded within its training set. Similarly, we might want a language model to \textit{be aware} of a misconception held by half of the population, but we certainly do not want the model to assert that misconception as true in 50\% of responses about it! Put differently, choosing the model's \emph{preferred responses and behaviors} from its broad \textit{knowledge and capabilities} is essential for constructing AI systems that are safe, effective, and controllable \citep{ouyang2022training}. Although current techniques usually guide LMs to align with human preferences through reinforcement learning (RL), we demonstrate that the RL-based objective commonly employed can be exactly optimized using a straightforward binary cross-entropy loss, substantially simplifying the preference learning process.

Broadly speaking, existing approaches embed desired behaviors into language models by utilizing carefully curated collections of human preferences that reflect behaviors considered safe and helpful. This preference learning phase follows an initial large-scale unsupervised pre-training on vast text corpora. The most direct method for preference learning is supervised fine-tuning on human demonstrations of high-quality outputs, but the most effective methods leverage reinforcement learning from human (or AI) feedback (RLHF/RLAIF; \citep{christiano2017deep,bai2022constitutional}). RLHF approaches involve fitting a reward model based on a dataset of human preferences and subsequently applying RL to optimize the language model policy to generate responses that yield high rewards while avoiding excessive deviation from the original model. Although RLHF yields models with impressive conversational and coding skills, the RLHF workflow is substantially more complicated than supervised learning, requiring training multiple language models and sampling from the policy during training, which results in considerable computational overhead.

In this work, we demonstrate a method to directly optimize a language model to align with human preferences without relying on explicit reward modeling or reinforcement learning techniques. We introduce \textit{Direct Preference Optimization} (DPO), an algorithm that implicitly maximizes the same objective as current RLHF approaches (reward maximization subject to a KL-divergence constraint) while being simpler to implement and easier to train. Conceptually, the DPO update increases the relative log probability of preferred responses compared to dispreferred ones, but it includes a dynamic, example-specific importance weight that prevents the model collapse observed when using a naive probability ratio objective. Similar to existing approaches, DPO is based on a theoretical preference model (such as the Bradley-Terry model; \cite{bradley1952rankanalysis}) that quantifies how well a reward function corresponds to empirical preference data. However, whereas current methods use this preference model to formulate a preference loss to train a reward model and then train a policy that optimizes this learned reward, DPO applies a change of variables to express the preference loss directly in terms of the policy. Consequently, given a dataset of human preferences over model outputs, DPO can optimize a policy through a straightforward binary cross entropy loss, yielding the optimal policy for an implicit reward function fitted to the preference data.

Our primary contribution is Direct Preference Optimization (DPO), a straightforward RL-free technique for training language models from preference data. Our experiments indicate that DPO performs at least as well as existing methods, including PPO-based RLHF, for preference learning in tasks like sentiment control, summarization, and dialogue, using language models with up to 6B parameters.

\section{Related Work}

Self-supervised language models of growing scale acquire the ability to perform certain tasks in a zero-shot manner \citep{radford2019language} or with few-shot prompting \citep{gpt3,megatron,chowdhery2022palm}. Nonetheless, their effectiveness on downstream tasks and alignment with user intentions can be greatly enhanced through fine-tuning on datasets containing instructions paired with human-generated completions \citep{mishra-etal-2022-cross,sanh2022multitask,chung2022scaling,thoppilan2022lamda}. This process, known as `instruction-tuning,' allows large language models (LLMs) to generalize to instructions beyond those seen during tuning and generally improves their practical usability \citep{chung2022scaling}. Although instruction tuning has been successful, collecting \textit{relative} human judgments about response quality is often simpler than obtaining expert demonstrations. Consequently, later studies have fine-tuned LLMs using datasets of human preferences, thereby enhancing capabilities in translation \citep{kreutzer-etal-2018-reliability}, summarization \citep{stiennon2022learning,ziegler2020finetuning}, storytelling \citep{ziegler2020finetuning}, and instruction-following \citep{ouyang2022training,ramamurthy2023is}. These approaches first train a neural network reward function to align with the preference dataset under models such as the Bradley-Terry model \citep{bradley1952rankanalysis}, then fine-tune the language model to maximize this reward using reinforcement learning techniques, frequently REINFORCE \citep{williams1992reinforce}, proximal policy optimization (PPO; \cite{schulman2017proximal}), or related variants \citep{ramamurthy2023is}. A related research direction employs LLMs fine-tuned for instruction-following with human feedback to produce additional synthetic preference data targeting specific attributes like safety or harmlessness \citep{bai2022constitutional}, relying solely on weak human supervision in the form of text rubrics for the LLM’s annotations. These approaches reflect a merging of two research strands: one focusing on training language models with reinforcement learning across various objectives \citep{Ranzato2015SequenceLT,paulus2018a,wu2018learning}, and another centered on general frameworks for learning from human preferences \citep{christiano2017deep,kupcsik2018learning}. Despite the advantages of utilizing relative human preferences, fine-tuning large language models with reinforcement learning remains a significant practical challenge; this study offers a theoretically grounded method for optimizing relative preferences without employing reinforcement learning.

Beyond the realm of language, the problem of learning policies from preferences has been explored in both bandit and reinforcement learning frameworks, with numerous methods put forward. Contextual bandit learning that relies on preferences or rankings of actions instead of explicit rewards is referred to as a contextual dueling bandit (CDB; \cite{yue2012karmed,dudik2015contextual}). When absolute rewards are unavailable, the theoretical treatment of CDBs replaces the concept of an optimal policy with a \textit{von Neumann winner}, defined as a policy whose expected victory rate against \textit{any} competing policy is at least 50\% \citep{dudik2015contextual}. Nonetheless, in the CDB framework, preference feedback is provided online, whereas in learning from human preferences, preference annotations are generally collected as a fixed offline batch of annotated action pairs \citep{yan2022human}. Similarly, \textit{preference-based RL} (PbRL) leverages binary preferences derived from an \textit{unknown} underlying `scoring' function in lieu of direct rewards \citep{BusaFekete2014,ruiz2023dueling}. A variety of PbRL algorithms exist, including ones capable of reusing off-policy preference data, but these typically involve two steps: first explicitly estimating the hidden scoring function (i.e., the reward model), and then optimizing the policy accordingly \citep{jain2013learning,BusaFekete2014,christiano2017deep,sadigh2017active,kupcsik2018learning}. In contrast, we propose a unified policy learning method that directly optimizes the policy to satisfy the observed preferences in a single stage.

\section{Preliminaries}\label{section:prelims}

We summarize the RLHF pipeline as described by \citeauthor{ziegler2020finetuning} and subsequent works \citep{stiennon2022learning, bai2022training, ouyang2022training}, which typically consists of three stages: 1) supervised fine-tuning (SFT); 2) preference sampling and reward modeling; and 3) reinforcement learning optimization.

\textbf{SFT}: The RLHF procedure generally starts by fine-tuning a pretrained language model using supervised learning on high-quality datasets tailored for the downstream tasks (such as dialogue or summarization), resulting in a model denoted by $\pi^\text{SFT}$.

\textbf{Reward Modeling Phase}: In this phase, the SFT model is prompted with inputs $x$ to generate pairs of responses $(y_1, y_2) \sim \pi^\text{SFT}(y \mid x)$. These pairs are then evaluated by human annotators who express a preference for one completion over the other, indicated as $y_w \succ y_l \mid x$, where $y_w$ and $y_l$ represent the preferred and less preferred outputs among $(y_1, y_2)$, respectively. It is assumed that these preferences arise from some hidden reward function $r^*(y, x)$ that we do not observe directly. Various methods exist to model these preferences, with the Bradley-Terry (BT) \cite{bradley1952rankanalysis} model being a commonly used approach (although more general Plackett-Luce ranking models \citep{plackett1975analysis, luce2012individual} can also be employed if multiple ranked responses are available). According to the BT model, the human preference distribution $p^*$ can be represented as:

\begin{equation}\label{eq:bradley-terry}
    p^*(y_1\succ y_2 \mid x)=\frac{\exp\left(r^*(x, y_1)\right)}{\exp\left(r^*(x, y_1)\right) + \exp\left(r^*(x, y_2)\right)}.
\end{equation}

Given access to a fixed dataset of comparisons $\mathcal{D}=\bigl\{x^{(i)}, y_w^{(i)}, y_l^{(i)}\bigr\}_{i=1}^N$ drawn from $p^*$, we can represent a reward model $r_{\phi}(x, y)$ and estimate its parameters through maximum likelihood. Treating this as a binary classification problem, the negative log-likelihood loss is expressed as:

\begin{equation}\label{eq:reward_model}
    \mathcal{L}_R(r_{\phi}, \mathcal{D}) = -\mathbb{E}_{(x, y_w, y_l)\sim \mathcal{D}}\bigl[\log \sigma(r_{\phi}(x, y_w)- r_{\phi}(x, y_l))\bigr]
\end{equation}

where $\sigma$ denotes the logistic function. In the setting of language models, $r_{\phi}(x, y)$ is frequently initialized from the SFT model $\pi^\text{SFT}(y \mid x)$ by appending a linear layer atop the final transformer layer to yield a single scalar reward prediction \cite{ziegler2020finetuning}. To reduce variance in the reward function, previous work normalizes rewards such that $\mathbb{E}_{x,y\sim \mathcal{D}}\left[r_\phi(x, y)\right] = 0$ for every $x$.

\textbf{RL Fine-Tuning Phase}: In the reinforcement learning stage, the trained reward function provides feedback to the language model. Building on previous approaches~\citep{jaques2017sequence, jaques2020human}, the optimization problem is formulated as

\begin{equation}\label{eq:RL}
\max_{\pi_{\theta}}  \mathbb{E}_{x\sim \mathcal{D}, y\sim \pi_{\theta}(y \mid x)}\bigl[r_{\phi}(x, y)\bigr] - \beta\mathbb{D}_{\textrm{KL}}\bigl[\pi_{\theta}(y\mid x)\mid \mid \pi_\text{ref}(y\mid x)\bigr],
\end{equation}

where $\beta$ governs how much the policy $\pi_{\theta}$ is allowed to diverge from the baseline reference policy $\pi_\text{ref}$, which is the initial SFT model $\pi^\text{SFT}$. In practical terms, the language model policy $\pi_\theta$ is initialized as $\pi^\text{SFT}$. This constraint is crucial as it limits the model’s departure from the distribution where the reward model provides accurate estimates, helps maintain diversity in generation, and prevents collapse to a few high-reward outputs. Because language generation is discrete, this objective is not differentiable and is generally optimized using reinforcement learning techniques. The common strategy \citep{ziegler2020finetuning, stiennon2022learning, bai2022training, ouyang2022training} defines the reward function as ${r(x, y) = r_{\phi}(x, y) -\beta (\log \pi_{\theta}(y\mid x) - \log \pi_\text{ref}(y\mid x))}$ and maximizes it using PPO \cite{schulman2017proximal}.

\section{Direct Preference Optimization}\label{sec:DPO}

Inspired by the difficulties of applying reinforcement learning methods to large-scale tasks like language model fine-tuning, we aim to develop a straightforward technique for policy optimization using preference data directly. Unlike earlier RLHF techniques that first learn a reward function and then optimize it via RL, our method uses a specific parameterization of the reward model that allows direct extraction of its optimal policy in closed form, eliminating the need for an RL training loop. As we explain in detail below, the central insight is to utilize an analytical relationship from reward functions to optimal policies, enabling the transformation of a loss over reward functions into a loss over policies. This change-of-variables strategy circumvents the necessity of training an explicit, separate reward model, while still optimizing with respect to established human preference models such as the Bradley-Terry model. Essentially, the policy network simultaneously encodes both the language model and the (implicit) reward.

\textbf{Derivation of the DPO objective.} We begin with the same reinforcement learning objective as previous studies, Eq.~\ref{eq:RL}, defined under a general reward function $r$. Following earlier work~\citep{peters2007reinforcement, peng2019advantage, korbak2022reinforcement, go2023aligning}, it is straightforward to demonstrate that the optimal solution to the KL-constrained reward maximization problem in Eq.~\ref{eq:RL} is given by:

\begin{equation}\label{eq:op_policy}
    \pi_r(y\mid x) = \frac{1}{Z(x)}\pi_\text{ref}(y\mid x)\exp\left(\frac{1}{\beta}r(x, y)\right),
\end{equation}

where $Z(x) =\sum_{y}\pi_\text{ref}(y\mid x)\exp\left(\frac{1}{\beta}r(x, y)\right)$ denotes the partition function. A full derivation is provided in Appendix \ref{app:derivation1}. Even when using the maximum likelihood estimate $r_{\phi}$ of the true reward function $r^*$, calculating the partition function $Z(x)$ remains computationally demanding \citep{korbak2022reinforcement, go2023aligning}, making this formulation impractical to employ directly. However, we can rearrange Eq.~\ref{eq:op_policy} to represent the reward function in terms of the corresponding optimal policy $\pi_r$, the reference policy $\pi_\text{ref}$, and the unknown partition function $Z(\cdot)$. Specifically, by taking the logarithm of both sides of Eq.~\ref{eq:op_policy} and performing some algebraic manipulation, we arrive at:

\begin{equation}\label{eq:main_eq}
    r(x,y) =\beta \log \frac{\pi_r(y\mid x)}{\pi_\text{ref}(y\mid x)} + \beta \log Z(x).
\end{equation}

This reparameterization can be applied to the true reward $r^*$ and its corresponding optimal policy $\pi^*$. Importantly, the Bradley-Terry model depends solely on the difference in rewards between two completions, i.e., ${p^*(y_1 \succ y_2 \mid x) = \sigma(r^*(x, y_1) - r^*(x, y_2))}$. Substituting the reparameterized form from Eq.~\ref{eq:main_eq} for $r^*(x,y)$ into the preference model Eq.~\ref{eq:bradley-terry}, the partition function cancels out, allowing the human preference probability to be expressed purely in terms of the optimal policy $\pi^*$ and reference policy $\pi_\text{ref}$. Consequently, the optimal RLHF policy $\pi^*$ under the Bradley-Terry framework fulfills the preference model:

\begin{equation}\label{eq:objective}
    p^*(y_1\succ y_2 \mid x)=\frac{1}{1 + \exp\left(\beta \log \frac{\pi^*(y_2\mid x)}{\pi_\text{ref}(y_2\mid x)} - \beta \log \frac{\pi^*(y_1\mid x)}{\pi_\text{ref}(y_1\mid x)}\right)}.
\end{equation}

The full derivation is detailed in Appendix~\ref{app:derivation2}. While Eq.~\ref{eq:objective} is based on the Bradley-Terry model, analogous formulations can be derived for the more general Plackett-Luce models~\citep{plackett1975analysis, luce2012individual}, as presented in Appendix~\ref{app:plackett_luce_models}.

Having expressed the probability of human preference data in terms of the optimal policy rather than the reward model, we can define a maximum likelihood objective for a parameterized policy $\pi_\theta$. Similar to the reward modeling approach (see Eq.~\ref{eq:reward_model}), the policy objective becomes:

\begin{equation}\label{eq:optimum_model}
    \mathcal{L}_\text{DPO}(\pi_{\theta}; \pi_\text{ref}) = -\mathbb{E}_{(x, y_w, y_l)\sim \mathcal{D}}\left[\log \sigma \left(\beta \log \frac{\pi_{\theta}(y_w\mid x)}{\pi_\text{ref}(y_w\mid x)} - \beta \log \

In this manner, we fit an implicit reward function using a different parameterization, whose optimal policy corresponds directly to $\pi_\theta$. Additionally, since our approach is equivalent to fitting a reparameterized Bradley-Terry model, it inherits certain theoretical guarantees, such as consistency under appropriate assumptions about the distribution of preference data \cite{bong2022generalized}. We further explore the theoretical aspects of DPO in comparison to other methods in Section~\ref{sec:theory}.

\textbf{What is the effect of the DPO update?} To gain a mechanistic insight into DPO, it is helpful to examine the gradient of the loss function $\mathcal{L}_\text{DPO}$. The gradient with respect to the parameters $\theta$ can be expressed as:

\begin{multline*}\label{eq:gradient}
    \nabla_\theta \mathcal{L}_\text{DPO}(\pi_\theta;\pi_\text{ref}) = \\ -\beta\mathbb{E}_{(x, y_w, y_l) \sim \mathcal{D}} \bigg[\underbrace{\sigma(\hat{r}_\theta(x, y_l) - \hat{r}_\theta (x, y_w))}_\text{larger weight when reward prediction is incorrect}\bigg[\underbrace{\nabla_\theta\log \pi(y_w \mid x)}_\text{boost likelihood of $y_w$} - \underbrace{\nabla_\theta\log\pi(y_l \mid x)}_\text{reduce likelihood of $y_l$}\bigg]\bigg],
\end{multline*}

where $\hat{r}_\theta(x, y) = \beta \log \frac{\pi_\theta(y \mid x)}{\pi_\text{ref}(y \mid x)}$ denotes the reward implicitly defined by the language model $\pi_\theta$ in relation to the reference model $\pi_\text{ref}$ (see Section~\ref{sec:theory} for details). Conceptually, the gradient of $\mathcal{L}_\text{DPO}$ increases the probability assigned to preferred outputs $y_w$ while decreasing it for less preferred outputs $y_l$. Crucially, the examples are weighted by the extent to which the implicit reward model $\hat{r}_\theta$ incorrectly favors the less preferred completions, scaled by $\beta$, reflecting the influence of the KL constraint. Our empirical results emphasize the significance of this weighting factor, as a naive version of this approach without it can lead to model collapse (see Appendix Table~\ref{tab:unlikelihood_generations}).

\textbf{DPO overview.} The standard DPO procedure includes the following steps: 1) For each prompt $x$, generate completions $y_1, y_2 \sim \pi_\text{ref}(\cdot \mid x)$ and annotate them with human preferences to form an offline preference dataset $\mathcal{D} = \{x^{(i)}, y_w^{(i)}, y_l^{(i)}\}_{i=1}^N$; 2) train the language model $\pi_\theta$ by minimizing the loss $\mathcal{L}_\text{DPO}$ given the reference policy $\pi_\text{ref}$, dataset $\mathcal{D}$, and target parameter $\beta$. 

In practical scenarios, it is preferable to leverage publicly accessible preference datasets instead of generating new samples and collecting human annotations. Because these preference datasets are generated under $\pi^\text{SFT}$, we set $\pi_\text{ref} = \pi^\text{SFT}$ whenever this model is available. If $\pi^\text{SFT}$ is unavailable, we initialize $\pi_\text{ref}$ by maximizing the likelihood of the preferred completions ${(x, y_w)}$, that is, by setting ${\pi_\text{ref} = \argmax_{\pi}\mathbb{E}_{x, y_w \sim \mathcal{D}}\left[\log \pi(y_w \mid x)\right]}$. This approach helps to reduce distributional mismatch between the true, unknown reference distribution and the $\pi_\text{ref}$ employed in DPO. Additional implementation specifics and hyperparameter details are provided in Appendix~\ref{app:implementation}.

\section{Theoretical Analysis of DPO} In this section, we further interpret the DPO approach, offer theoretical justification, and discuss how DPO's benefits address challenges found in actor-critic algorithms commonly used for RLHF, such as PPO~\cite{schulman2017proximal}.

\label{sec:theory}

\subsection{Your Language Model Is Secretly a Reward Model} DPO circumvents the need to explicitly fit a reward function and avoids reinforcement learning by optimizing a single maximum likelihood objective. Notice that the optimization problem in Eq. \ref{eq:main_eq} is equivalent to a Bradley-Terry model where the reward is parameterized as $r^*(x, y) = \beta \log\frac{\pi^*_\theta(y \mid x)}{\pi_\text{ref}(y \mid x)}$. We thereby train the parametric model $\pi_{\theta}$ in a manner analogous to the reward model optimization in Eq. \ref{eq:reward_model}, but after applying a change of variables. In this section, we develop the theoretical foundation behind this reparameterization, demonstrate that it does not restrict the family of learnable reward models, and prove it enables exact recovery of the optimal policy. We start by defining an equivalence relation between reward functions.

\begin{definition}
We define two reward functions $r(x, y)$ and $r'(x, y)$ as equivalent if and only if there exists some function $f$ such that ${r(x, y)-r'(x, y) = f(x)}$.     
\end{definition}

It is straightforward to verify that this indeed constitutes an equivalence relation, dividing the set of reward functions into distinct equivalence classes. We can then present the following two lemmas:

\begin{lemma}\label{lemma:same_prefrence} Within the Plackett-Luce framework, and specifically the Bradley-Terry model, any two reward functions belonging to the same equivalence class generate the identical preference distribution.
\end{lemma}

\begin{lemma}\label{lemma:same_policy}
    Any two reward functions from the same equivalence class produce the same optimal policy for the constrained RL problem.
\end{lemma}

The proofs are direct and are provided in Appendix \ref{app:lemma1}. The first lemma highlights a well-known under-specification issue characteristic of the Plackett-Luce family of models \cite{plackett1975analysis}. Because of this under-specification, it is typically necessary to impose extra identifiability constraints to obtain guarantees on the MLE estimates derived from Eq. \ref{eq:reward_model} \cite{bong2022generalized}. The second lemma indicates that all reward functions within the same class correspond to the same optimal policy; therefore, when optimizing the final objective, it suffices to recover any reward function from this optimal equivalence class. We establish the following theorem in Appendix~\ref{app:thm1}:

\begin{theorem}\label{thm:main}
    Under mild conditions, every reward equivalence class that aligns with the Plackett-Luce (particularly Bradley-Terry) models can be expressed through the reparameterization ${r(x, y) = \beta \log \frac{\pi(y\mid x)}{\pi_\text{ref}(y\mid x)}}$, for some model $\pi(y\mid x)$ and a given reference model $\pi_\text{ref}(y \mid x)$.
\end{theorem}

\begin{sproof}
    Take any reward function $r(x, y)$, which defines an associated optimal model $\pi_r(y \mid x)$ as given by Eq. \ref{eq:op_policy}. We will demonstrate that a reward function within the equivalence class of $r$ admits the reparameterization stated above. Define the projection $f$ by  
\begin{equation}
    f(r; \pi_\text{ref}, \beta)(x, y) = r(x, y) - \beta\log\sum_{y}\pi_\text{ref}(y\mid x)\exp\left(\frac{1}{\beta}r(x, y)\right)
\end{equation}
The operator $f$ effectively normalizes the reward function by subtracting the logarithm of the partition function corresponding to $\pi_r$. Since this normalization term depends solely on the prefix $x$, the resulting function $f(r; \pi_\text{ref}, \beta)(x, y)$ lies within the equivalence class of $r(x, y)$. Finally, substituting $r$ with the expression on the right-hand side of Eq.~\ref{eq:main_eq} (valid for any reward function), yields $f(r; \pi_\text{ref}, \beta)(x, y) = \beta \log \frac{\pi_r(y\mid x)}{\pi_\text{ref}(y\mid x)}$. Hence, the projection $f$ constructs an element of $r$’s equivalence class in the desired form, showing that the proposed reparameterization of the reward function does not limit generality.
\end{sproof}

Alternatively, Theorem~\ref{thm:main} can be interpreted as identifying the specific reward function chosen by the DPO reparameterization within each equivalence class, namely, the reward function that fulfills the condition:

\begin{equation}\label{eq:lag_p}
     \sum_{y}\underbrace{\pi_\text{ref}(y\mid x)\exp\left(\frac{1}{\beta}r(x, y)\right)}_{=\pi(y\mid x)\text{, using Thm.~\ref{thm:main} reparam.}} = 1,
\end{equation}

which means that $\pi(y\mid x)$ constitutes a valid probability distribution (i.e., all probabilities are positive and their sum equals 1). From Eq.~\ref{eq:op_policy}, it is evident that Eq.~\ref{eq:lag_p} represents the partition function of the optimal policy derived from the reward function $r(x, y)$. The crucial insight underlying the DPO algorithm is that by enforcing certain constraints on the otherwise underdetermined Plackett-Luce family of preference models (specifically the Bradley-Terry model), we retain the entire class of representable reward functions while ensuring that the optimal policy described in Eq.~\ref{eq:op_policy} is analytically computable for every prompt $x$.

\subsection{Instability of Actor-Critic Algorithms} Our framework also enables us to analyze the instability issues encountered with conventional actor-critic methods employed in RLHF, such as PPO. Following the RLHF pipeline, we concentrate on the reinforcement learning fine-tuning step detailed in Section~\ref{section:prelims}. We establish links to the control-as-inference paradigm \cite{levine2018reinforcement} in the context of the constrained RL problem described in Eq.~\ref{eq:RL}. Suppose we have a parameterized policy model $\pi_{\theta}(y\mid x)$ and seek to minimize the KL divergence $\mathbb{D}_{\text{KL}}[\pi_{\theta}(y|x) \mid \mid \pi^*(y\mid x)]$, where $\pi^*$ denotes the optimal policy from Eq.~\ref{eq:optimum_model} induced by the reward function $r_{\phi}(y, x)$. Through straightforward algebraic manipulation, this yields the optimization objective:

\begin{equation}\label{eq:AC}
    \max_{\pi_{\theta}}\mathbb{E}_{\pi_{\theta}(y\mid x)}\bigg[\underbrace{r_{\phi}(x, y) -\beta\log\sum_{y}\pi_\text{ref}(y\mid x)\exp\left(\frac{1}{\beta}r_{\phi}(x, y)\right)}_{f(r_{\phi}, \pi_\text{ref}, \beta)} - \underbrace{\beta\log\frac{\pi_{\theta}(y\mid x)}{\pi_\text{ref}(y\mid x)}}_{\text{KL}}\bigg].
\end{equation}

This objective is identical to the one optimized in previous studies \citep{ziegler2020finetuning, stiennon2022learning, bai2022training, ouyang2022training} that use the DPO-equivalent reward for the reward class $r_{\phi}$. In this framework, the normalization component in $f(r_{\phi}, \pi_\text{ref}, \beta)$ can be viewed as the soft value function associated with the reference policy $\pi_\text{ref}$. Although this term does not influence the optimal solution, omitting it can lead to high variance in the policy gradient of the objective, which may cause unstable learning. We can account for the normalization term by employing a learned value function, but this approach can be challenging to optimize. Alternatively, earlier works have normalized rewards by using a human completion baseline, effectively a single-sample Monte-Carlo approximation of the normalization factor. In contrast, the DPO reparameterization produces a reward function that does not require any baselines.

\section{Experiments}
This section presents an empirical evaluation of DPO's capability to train policies directly from preference data. Initially, in a controlled text-generation environment, we investigate: how effectively does DPO balance maximizing reward while minimizing the KL-divergence with the reference policy compared to standard preference learning methods like PPO? Subsequently, we assess DPO's performance on larger models and more complex RLHF tasks, including summarization and dialogue. Our findings indicate that, with minimal hyperparameter tuning, DPO typically matches or surpasses strong baselines such as RLHF with PPO and also outperforms methods that select the best of $N$ sampled trajectories under a learned reward function. Before detailing these results, we outline the experimental setup; further information is provided in Appendix~\ref{app:exp_details}.

\textbf{Tasks.} Our study investigates three distinct open-ended text generation tasks. Across all experiments, the algorithms learn a policy based on a preference dataset $\mathcal{D}=\bigl\{x^{(i)}, y_w^{(i)}, y_l^{(i)}\bigr\}_{i=1}^N$. In the case of \textbf{controlled sentiment generation}, $x$ represents a prefix of a movie review from the IMDb dataset \cite{maas-EtAl:2011:ACL-HLT2011}, and the policy is tasked with producing $y$ that conveys a positive sentiment. To facilitate a controlled evaluation, we \textit{generate} preference pairs from the generated outputs using a pre-trained sentiment classifier, where the probability $p(\text{positive}\mid x,y_w)$ is greater than $p(\text{positive}\mid x,y_l)$. For SFT, GPT-2-large is fine-tuned until convergence on reviews from the IMDb training split (additional details in App~\ref{app:sentiment_details}). In the \textbf{summarization} task, $x$ is a Reddit forum post, and the policy must generate a summary $y$ that captures the key points of the post. Following previous work, we use the Reddit TL;DR summarization dataset \citep{volske-etal-2017-tl} alongside human preferences collected by \citeauthor{stiennon2022learning}. We employ an SFT model fine-tuned on human-written summaries of forum posts\footnote{\url{https://huggingface.co/CarperAI/openai_summarize_tldr_sft}} using the TRLX \citep{leandro_von_werra_2023_7790115} framework for RLHF. The human preference dataset was obtained by \citeauthor{stiennon2022learning} from samples generated by a different, though similarly fine-tuned, SFT model. Lastly, in the \textbf{single-turn dialogue} task, $x$ corresponds to a user query, which could range from astrophysics questions to requests for relationship advice. The policy must generate an engaging and useful response $y$ to the user's input. We utilize the Anthropic Helpful and Harmless dialogue dataset \citep{bai2022training}, which includes 170k dialogues between humans and an automated assistant. Each transcript concludes with a pair of responses produced by a large (but unspecified) language model, accompanied by a preference label indicating the response favored by humans. Since no pre-trained SFT model is available in this context, we fine-tune a standard language model solely on the preferred completions to create the SFT model.

\textbf{Evaluation.} Our experiments employ two distinct evaluation strategies. To assess how effectively each algorithm optimizes the constrained reward maximization objective, in the controlled sentiment generation setting we measure each algorithm by the frontier of achieved reward and KL-divergence relative to the reference policy; this frontier is calculable since the ground-truth reward function (a sentiment classifier) is available. However, in practical scenarios, the true reward function is unknown; hence, we assess algorithms based on their \textit{win rate} against a baseline policy, utilizing GPT-4 as a stand-in for human assessment of summary quality and response helpfulness in the summarization and single-turn dialogue contexts, respectively. For summarization, the baseline is the reference summaries in the test set; for dialogue, it is the preferred response from the test dataset. While prior research indicates that LMs may serve as superior automated evaluators compared to existing metrics \citep{Chen2023ExploringTU}, we perform a human study to validate our use of GPT-4 for evaluation in Sec.~\ref{sec:human-judgments}. We observe that GPT-4 evaluations strongly correlate with human judgments, with human agreement with GPT-4 generally matching or exceeding the agreement among human annotators.

\textbf{Methods.} Besides DPO, we test several established methods for training language models to align with human preferences. At a basic level, we investigate zero-shot prompting with \textbf{GPT-J} \citep{gpt-j} for the summarization task and 2-shot prompting with \textbf{Pythia-2.8B} \citep{biderman2023pythia} for the dialogue task. We also evaluate the \textbf{SFT} model and \textbf{Preferred-FT}, which is fine-tuned via supervised learning on the selected completion $y_w$ derived either from the SFT model (in controlled sentiment and summarization) or from a generic LM (in single-turn dialogue). Another pseudo-supervised approach is \textbf{Unlikelihood}~\citep{welleck2019neural}, which optimizes the policy to increase the likelihood of $y_w$ while \textit{decreasing} the likelihood of $y_l$; an optional weighting coefficient $\alpha \in [0,1]$ is applied to the unlikelihood component. We additionally consider \textbf{PPO} \citep{schulman2017proximal}, which uses a reward function trained from preference data, as well as \textbf{PPO-GT}, an oracle method that learns from the ground truth reward function available in the controlled sentiment setting. For sentiment experiments, we employ two implementations of PPO-GT: an off-the-shelf version \cite{leandro_von_werra_2023_7790115} and a modified variant that normalizes rewards and tunes hyperparameters for improved performance (these modifications are also applied to the standard PPO using learned rewards). Lastly, we include the \textbf{Best of $N$} baseline, which samples $N$ responses from the SFT model (or Preferred-FT for dialogue) and selects the highest-scoring response based on a reward model trained on preference data. Although this effective method separates reward model quality from PPO optimization, it is computationally costly even for moderate $N$, as it requires generating $N$ completions per query during evaluation.

\subsection{How effectively can DPO optimize the RLHF objective?}

The KL-constrained reward maximization objective employed in standard RLHF algorithms balances maximizing reward while limiting the policy's deviation from the reference policy. Consequently, when evaluating algorithms, it is important to consider both the obtained reward and the KL divergence; obtaining a slightly higher reward at the cost of a significantly increased KL divergence may not be favorable. Figure~\ref{fig:frontier-tldr-main} illustrates the reward-KL frontier for different algorithms within the sentiment task. We conduct multiple training experiments for each algorithm, varying the hyperparameter controlling policy conservativeness in each run (target KL values $\in \{3,6,9,12\}$ for PPO, $\beta \in \{0.05,0.1,1,5\}$, $\alpha \in \{0.05,0.1,0.5,1\}$ for unlikelihood, and random seeds for preferred-FT). This hyperparameter sweep comprises 22 runs in total. After every 100 training steps until convergence, we evaluate each policy on a set of test prompts, calculating the average reward under the true reward function along with the average sequence-level KL\footnote{Specifically, the sum of the per-timestep KL divergences.} relative to the reference policy, $\text{KL}\left(\pi \mid\mid \pi_\text{ref}\right)$. Our findings indicate that DPO establishes the most efficient frontier by a large margin, attaining the highest reward while maintaining a low KL divergence. This outcome is particularly significant for several reasons. First, although DPO and PPO optimize the same objective, DPO is distinctly more efficient; the reward/KL tradeoff achieved by DPO strictly surpasses that of PPO. Second, DPO outperforms PPO even when PPO has access to ground truth rewards (PPO-GT).

\subsection{Is DPO scalable to actual preference datasets?}
\label{sec:dpo-real-datasets}
We proceed to assess the fine-tuning effectiveness of DPO on tasks involving summarization and single-turn dialogue. In the context of summarization, automated metrics like ROUGE often show weak alignment with human judgments~\citep{stiennon2022learning}, and previous research has demonstrated that fine-tuning language models with PPO using human preference data results in more effective summaries. Our evaluation compares different methods by generating completions on the test portion of the TL;DR summarization dataset and calculating the average win rate against the reference completions from the test set. For all approaches, completions are sampled across a temperature range from 0.0 to 1.0, with the win rates displayed in Figure~\ref{fig:frontier-tldr-main} (right). DPO, PPO, and Preferred-FT all fine-tune the identical GPT-J SFT model\footnote{\url{https://huggingface.co/CarperAI/openai_summarize_tldr_sft}}. Our findings indicate that DPO attains a win rate around 61\% at temperature 0.0, outperforming PPO which achieves approximately 57\% at its optimal sampling temperature of 0.0. Furthermore, DPO reaches a higher peak win rate compared to the best-performing baseline among $N$ methods. It is worth noting that we did not extensively optimize DPO's $\beta$ hyperparameter, suggesting these results might underestimate its full capability. Additionally, DPO exhibits significantly greater resilience to changes in sampling temperature relative to PPO, whose performance can decline to that of the underlying GPT-J model at elevated temperatures. Preferred-FT shows no notable improvement over the SFT baseline. We also perform a direct comparison between DPO and PPO in human evaluations discussed in Section~\ref{sec:human-judgments}, where samples generated by DPO at temperature 0.25 were favored 58\% of the time over PPO samples at temperature 0.

For single-turn dialogue, we assess various methods using a subset of the test split from the Anthropic HH dataset \citep{bai2022training}, which involves one round of human-assistant interaction. GPT-4 evaluations utilize the preferred completions from the test set as a benchmark to calculate the win rate across different methods. Since there is no established SFT model for this task, we begin with a pre-trained Pythia-2.8B, apply Preferred-FT to train a reference model on the selected completions to ensure the outputs remain within the model’s distribution, and subsequently train using DPO. We also compare against the best among 128 Preferred-FT completions (noting that the Best of $N$ baseline plateaus at 128 completions for this task; see Appendix Figure~\ref{fig:best-of-n}) and a 2-shot prompted variant of the Pythia-2.8B base model, finding that DPO performs equally well or better at optimal temperatures for each method. Additionally, we evaluate an RLHF model trained via PPO on the Anthropic HH dataset\footnote{\url{https://huggingface.co/reciprocate/ppo_hh_pythia-6B}} from a recognized source\footnote{\url{https://github.com/CarperAI/trlx/tree/main/examples/hh}}, but we are unable to identify a prompt or sampling temperature that surpasses the baseline Pythia-2.8B model's performance. Drawing from our TL;DR results and the fact that both methods optimize the same reward, we treat Best of 128 as an approximate benchmark for PPO-level performance. In summary, DPO stands out as the only computationally efficient approach that outperforms the preferred completions in the Anthropic HH dataset and achieves comparable or superior results relative to the computationally intensive Best of 128 baseline. Lastly, Figure~\ref{fig:dialogue-main} demonstrates that DPO reaches its peak performance relatively rapidly.

\subsection{Generalization to a new input distribution}

To further assess the performance differences between PPO and DPO under shifts in data distribution, we tested the PPO and DPO policies from our Reddit TL;DR summarization experiment on a distinct dataset: news articles from the test split of the CNN/DailyMail dataset \citep{nallapati-etal-2016-abstractive}. We employed the optimal sampling temperatures identified in the TL;DR task (0 and 0.25). The outcomes are displayed in Table~\ref{tab:ood}. We calculated GPT-4's win rate compared to the ground-truth summaries in the datasets, using the same GPT-4 (C) prompt from the Reddit TL;DR experiment, substituting the phrase ``forum post'' with ``news article''. For this new dataset, DPO consistently surpasses the PPO policy by a notable margin. This experiment offers preliminary evidence that DPO policies are capable of generalizing at least as effectively as PPO policies, despite DPO not leveraging the additional unlabeled Reddit TL;DR prompts that PPO utilizes.

\subsection{Validating GPT-4 judgments with human judgments}
\label{sec:human-judgments}

We carried out a human evaluation to confirm the trustworthiness of GPT-4’s assessments, using results from the TL;DR summarization experiment with two distinct GPT-4 prompts. The \textbf{GPT-4 (S)} (simple) prompt requests a choice of which summary better captures the important information from the post. The \textbf{GPT-4 (C)} (concise) prompt asks additionally which summary is more concise; we included this prompt because we observed that GPT-4 tends to favor longer and more repetitive summaries than humans do when using the \textbf{GPT-4 (S)} prompt. The full prompts are provided in Appendix~\ref{app:prompts}. We performed three pairwise comparisons involving the highest-performing method (DPO, temp. 0.25), the lowest-performing method (PPO, temp. 1.0), and a mid-level method (SFT, temp. 0.25), aiming to capture a range of summary qualities. Each of these methods was compared against greedily-sampled PPO (at its best-performing temperature). Our findings indicate that for both prompts, GPT-4’s agreement with human judgments occurs about as frequently as humans agree among themselves, implying that GPT-4 can serve as a reliable stand-in for human evaluation (due to a limited number of human raters, multiple human judgments were collected only for the DPO vs. PPO-1 comparison). Overall, the \textbf{GPT-4 (C)} prompt yields win rates that align more closely with human preferences; consequently, this prompt is used for the main results reported in Section~\ref{sec:dpo-real-datasets}. Additional information regarding the human evaluation study, including the rater interface and a list of volunteers, can be found in Appendix~\ref{app:human-study}.

\section{Discussion}
Learning from preferences offers a robust and scalable approach for developing capable and aligned language models. We have presented DPO, a straightforward training method that enables language models to be trained from preferences without the need for reinforcement learning. Instead of forcing the preference learning problem into a conventional RL framework to utilize existing RL algorithms, DPO establishes a connection between language model policies and reward functions, allowing for the direct training of a language model to meet human preferences \textit{directly} using a simple cross-entropy loss, eliminating the necessity for reinforcement learning and preserving generality. With minimal hyperparameter tuning, DPO achieves performance comparable to or surpassing current RLHF techniques, including those based on PPO, thereby significantly lowering the barrier to training more language models using human preferences.

\textbf{Limitations \& Future Work.} Our findings prompt several critical questions for subsequent investigation. How well does the DPO policy generalize outside the training distribution compared to approaches relying on an explicit reward function? Preliminary results indicate that DPO policies may generalize similarly to PPO-based models, but a more thorough examination is required. For instance, can self-labeling using the DPO policy enable effective utilization of unlabeled prompts? Another avenue to explore is how reward over-optimization appears in the context of direct preference optimization and whether the modest performance decline observed in Figure~\ref{fig:dialogue-main}-right reflects this phenomenon. Moreover, although our evaluations cover models up to 6B parameters, scaling DPO to state-of-the-art models that are magnitudes larger represents an exciting direction for future research. Concerning evaluation, we observe that GPT-4’s computed win rates are influenced by the prompt, suggesting that future work could investigate optimal methods for eliciting high-quality assessments from automated evaluators. Lastly, DPO has numerous potential applications beyond training language models from human preferences, such as training generative models across other modalities.

\end{document}